% ---------------------------------------------------------------------
%  第 13 章 软件工程
% ---------------------------------------------------------------------
\chapter{软件工程}
\label{ch:engineering}

前十二章从数值与物理角度描述了库。本章换一个视角：把它当作一个
软件项目，说明目录组织、模块边界、可见性策略、文档规范与规模。
这些选择直接影响了代码的可维护性与可审计性。

\section{目录组织}
\label{sec:layout}

\cref{fig:src-tree}列出源码模块及其职责，各源文件的行数可查阅\cref{tab:file-index}。

\begin{figure}[htbp]
    \centering
    \begin{tikzpicture}[y=.56cm, every node/.style={font=\small}]
        \node[draw=kbblue!60,fill=kbblue!8,rounded corners=2pt,anchor=west]
            (root) at (0,-7.5) {\code{src/}};
        \draw[kbblue!60] (root.east) -- (1.8,-7.5) (1.8,0) -- (1.8,-15);
        \foreach \entry/\desc [count=\row from 0] in {
            base/{Dyson/VIE2 共享求解内核},
            storage/{容器存储访问},
            algebra/{容器代数},
            variants/{容器家族操作},
            complex/{复数类型与算术},
            convolution/{轮廓卷积与气泡图},
            dyson/{Dyson 方程求解},
            equilibrium/{平衡态与稳态构造},
            fourier/{傅里叶变换},
            io/{文本输入与输出},
            linalg/{矩阵运算与线性求解},
            math/{分布函数、插值与外推},
            observables/{密度、能量、电流与谱函数},
            types/{轮廓与格林函数容器},
            vie2/{第二类 Volterra 积分方程求解},
            weights/{积分与差分权重}
        } {
            \draw[kbblue!60] (1.8,-\row) -- (2.25,-\row);
            \node[anchor=west] at (2.3,-\row) {\texttt{\entry}};
            \node[anchor=west] at (5.15,-\row) {\desc};
        }
    \end{tikzpicture}
    \caption[源码目录]{源码目录。右侧逐行列出各模块及其主要职责。}
    \label{fig:src-tree}
\end{figure}

目录按数值职责划分。例如 \code{convolution} 使用容器、权重与线性代数，
\code{dyson} 与 \code{vie2} 则调用卷积内核。审查卷积改动时，可沿这条
调用关系检查推迟、左混合和小于分量的调用约定。

库源码由同一个 crate 编译，目录边界主要用来控制依赖和定位问题。
卷积内核的索引改动可先用 \code{cargo test --test convolution} 检查分量级结果，
再运行 Dyson 与 VIE2 的集成测试，检查同一内核在启动与推进中的调用。
\code{tests/} 负责按主题组织这些测试，\code{examples/} 则给出从配置到输出的完整模型驱动。

\section{模块边界}
\label{sec:module-boundary}

容器层定义分量的形状与索引，积分权重层提供与模型无关的求积系数。
卷积将两者组合起来。它读取输入分量块，并将带权乘积累加到调用者提供的目标。
求解器在此基础上确定当前时间步的计算顺序，并负责把卷积源项和隐式系数组装进块系统。
平衡态构造与 Fourier 变换还会为求解路径提供初值和频域数据。

\paragraph{私有内核的下沉。}
只在一个模块中使用的辅助函数可以保持私有。卷积和求解器之间共享的内核，
可通过 \code{pub(crate)} 提供给同一 crate 使用。公共入口的可达性由其所在模块
及重导出路径共同决定。查找入口时先看模块的 \code{mod.rs}，再进入具体实现文件。

\paragraph{工作区的所有权。}
卷积工作区保存随块尺寸确定的中间乘积。完整轮廓驱动还可以在整个实时间循环中
复用虚时乘积缓存。求解器的块工作区保存组装矩阵、右端与枢轴数据。
同一系数矩阵的多个右端可以共享一次分解。系数随时间变化时，仍需重新分解，
但组装和求解使用的存储可以复用。这样的生命周期区分便于判断某一处优化减少的是分配、
重复乘积，还是分解次数。

\section{可见性策略}
\label{sec:visibility}

\cref{tab:visibility}概括不同可见性声明在库中的用途。

库的可见性分为三级：
\begin{itemize}
    \item \code{pub}：公共应用程序编程接口（Application Programming Interface，API）的条目在模块的 \code{mod.rs} 中重新导出；
    \item \code{pub(crate)}：跨模块内部使用，不出现在 \code{cargo doc}
          的公共页面中；
    \item 私有：只在定义它的模块内可见。
\end{itemize}

\begin{table}[htbp]
    \centering
    \caption[可见性与规模]{可见性与规模}
    \label{tab:visibility}
    \small
    \begin{tabularx}{\textwidth}{@{}lX@{}}
        \toprule
        类别 & 说明 \\
        \midrule
        \code{pub fn} & 约 550 处，包含各模块的公共入口与少量容器方法 \\
        \code{pub(crate) fn} & 约 120 处，主要是内核、检查函数与共享体 \\
        私有函数 & 其余；占比最高，多为单一调用点的辅助 \\
        \bottomrule
    \end{tabularx}
\end{table}

\code{mod.rs} 中的导出按可见性分组：先是 \code{pub use}，
再是 \code{pub(crate) use}。这一分组让读者一眼看出哪些是公共
接口，哪些是内部约定。

\section{文档规范}
\label{sec:doc-style}

模块文档用 \code{//!} 说明分量约定和算法背景，函数文档用 \code{///}
说明入口的参数与行为。对数值接口，最重要的是明确时间指标范围、
矩阵作用方向、统计符号，以及输出是覆盖还是累加。
以小于分量的增量卷积为例，目标按较晚时间组织，内核在已有目标上累加，
调用者因此需要按求解流程决定何时清零。离散公式中的指标方向也应与循环边界一致。

\paragraph{数值变更的验证。}
块乘法的实现可以先用随机矩阵与通用路径比较。积分项的重排还需要检查
启动区、积分端点、单位权重内部区以及伴随重构。分量级测试覆盖这些局部行为，
集成测试用参考数据检验多个分量组成的传播结果，阶数门则用网格加密检查误差的变化趋势。
三种证据回答不同问题，测试范围应与改动可能影响的路径对应。

\section{自实现清单}
\label{sec:self-implemented}

库不依赖任何 crate。为此自实现的部分包括：
\begin{itemize}
    \item 复数类型与全套算术；
    \item 稠密矩阵、LU 分解、块线性求解；
    \item 厄米特征值分解（循环 Jacobi）；
    \item 矩阵指数（缩放--平方 + Padé）；
    \item 任意长度 FFT（混合基数 + Bluestein）；
    \item 有理数类型（用于权重表的精确构造）；
    \item 文本解析与格式化（包括 C \code{\%g} 的复刻）；
    \item TOML（Tom's Obvious, Minimal Language）子集解析，供算例读取配置。
\end{itemize}

自实现的边界是“标准库能提供的不用自己写”：
\code{Vec}、\code{Arc}、\code{OnceLock}、\code{Mutex} 与文件读写
都来自标准库；数值算法则全部自写。这一边界让库的依赖图为空，
同时保留标准库提供的容器、同步与文件操作接口。算例配置解析与模型驱动位于支持层，
求解器接收已构造的容器和参数，可在不同的输入配置与输出流程中复用。

\section{命名约定}
\label{sec:naming}

\cref{tab:naming}给出区分分量、目标类型和内核变体的命名示例。

\begin{table}[htbp]
    \centering
    \caption[后缀与前缀约定]{后缀与前缀约定}
    \label{tab:naming}
    \small
    \begin{tabularx}{\textwidth}{@{}llX@{}}
        \toprule
        模式 & 含义 & 示例 \\
        \midrule
        \code{_mat} & Matsubara 分量上的操作 &
            \code{dyson_mat}, \code{conv_mat} \\
        \code{_ret} / \code{_less} / \code{_lmix} &
            特定分量上的操作 &
            \code{vie2_timestep_ret} \\
        \code{_step} & 目标是时间步对象 &
            \code{incr_convolution_ret_step} \\
        \code{_buf} & 目标是块缓冲 &
            \code{incr_less_adv_buf} \\
        \code{_impl} & 多个公共入口的共享体 &
            \code{incr_convolution_less_impl} \\
        \code{_d1} & 标量专用内核 &
            \code{conv_ret_kernel_d1} \\
        \code{_herm} & 利用厄米性减少一半计算 &
            \code{bubble1_herm} \\
        \code{_kms} & 利用虚时反射关系 &
            \code{conv_mat_kms} \\
        \code{_moving} & 移动窗版本 &
            \code{dyson_timestep_moving} \\
        \code{_pseudo} & 赝粒子版本 &
            \code{pseudo_dyson} \\
        \bottomrule
    \end{tabularx}
\end{table}

\section{规模统计}
\label{sec:size}

\cref{tab:module-size}按模块统计源文件、行数和函数声明。

文件数统计 \code{src/} 下的 Rust 文件。行数包括空行、注释和文件内的单元测试。
函数列统计 \code{pub fn} 与 \code{pub const fn} 声明，实际公共调用路径还取决于模块可见性和重导出。

\begin{table}[htbp]
    \centering
    \caption[各模块的规模]{各模块的规模}
    \label{tab:module-size}
    \small
    \begin{tabular}{@{}lrrr@{}}
        \toprule
        模块 & 文件数 & 行数 & \code{pub fn} 数 \\
        \midrule
% BEGIN GENERATED SIZE TABLE (gen_modules.py)
        \code{complex} & 4 &  387 & 8 \\
        \code{linalg} & 7 & 2180 & 28 \\
        \code{math} & 7 & 2592 & 22 \\
        \code{types} & 14 & 2700 & 99 \\
        \code{base} & 5 &  571 & 2 \\
        \code{storage} & 7 & 1274 & 50 \\
        \code{algebra} & 5 & 2755 & 87 \\
        \code{variants} & 4 & 1255 & 32 \\
        \code{weights} & 8 & 1338 & 13 \\
        \code{convolution} & 12 & 7194 & 43 \\
        \code{dyson} & 9 & 3581 & 28 \\
        \code{vie2} & 9 & 3264 & 24 \\
        \code{equilibrium} & 6 & 1852 & 8 \\
        \code{fourier} & 7 & 2724 & 26 \\
        \code{observables} & 5 &  502 & 12 \\
        \code{io} & 8 & 1989 & 11 \\
        \code{lib.rs} & 1 &   53 & 0 \\
        \midrule
        合计 & 118 & 36\,211 & 493 \\
% END GENERATED SIZE TABLE (gen_modules.py)
        \bottomrule
    \end{tabular}
\end{table}

卷积模块是最大的单一模块，约占总行数的五分之一；这与它的角色相符：
三类分量 × 四种调用形态 × 三种目标类型，组合数最多。其次是
\code{algebra}/\code{variants}（容器代数与家族操作）、\code{base}
（共享求解内核）与求解器两个模块。
嵌入源码的单元测试计入行数。独立的集成测试位于 \code{tests/}，其分布见附录\ref{app:tests}。
